Hand-crafted microphone-array geometry priors are widely injected into deep
sound event localization and detection (SELD) systems, often on the assumption
that explicit spatial layout can only help. We test this assumption with a
one-bit intervention in the official DCASE~2024 SELD baseline: a geometry bias is
toggled on or off inside an otherwise identical channel-attention block. Across
two input modalities (MIC and FOA), three temporal backbones (CRNN+MHSA,
Conformer, Transformer), and five paired seeds on real STARSS23 recordings, the
prior's effect on directional accuracy is not universal but graded by temporal
inductive bias. It helps the recurrent FOA model
($-5.0\degsym$ class-dependent localization error, $\dz=-1.03$), is near-neutral
in the Conformer, and harms the pure-attention Transformer (up to
$+5.7\degsym$, $\dz=+0.71$). A factorial test over the recurrent and Transformer
extremes confirms an architecture$\times$prior interaction ($F=4.62$,
$p=0.039$; seed-matched paired contrast $p=0.029$), and the recurrent FOA
benefit replicates zero-shot on STARSS22 ($\dz=-1.50$, $p=0.029$). The same
architecture axis does not govern distance: relative-distance gains are instead
concentrated in Ambisonics by effect size. A linear probe finds no significant
change in linearly decodable direction, indicating a behavioural rather than
simple representational shift. Finally, a second injection mechanism that adds
geometry as an early convolutional feature bias exposes the boundary of this
claim: FOA contrasts move under rerun and validation-selection changes, and one
MIC-Transformer contrast reverses relative to GCA.
We conclude that geometry priors are architecture-structured but
injection- and training-selection-dependent design choices, not universal SELD
improvements.
